\documentclass[10pt,a4paper]{amsart}
\usepackage[margin=19mm]{geometry}
\usepackage{amsmath,amssymb,mathtools}
\usepackage[scr=rsfs,cal=euler]{mathalfa}
\usepackage{xfrac}
\usepackage[unicode,hidelinks,bookmarks=false]{hyperref}
\newcommand{\ie}{i.e.\ }
\newcommand{\cf}{cf.\ }
\newcommand{\resp}{respectively\ }
\newcommand{\Subsection}{\subsection}
\providecommand{\nobreakdash}{\nobreak\hbox{-}}
\setlength{\emergencystretch}{2em}
\multlinegap0pt
\theoremstyle{plain}
\newtheorem{thm}{Theorem}
\newtheorem{lemma}[thm]{Lemma}
\newtheorem{cor}[thm]{Corollary}
\newtheorem{prop}[thm]{Proposition}
\newtheorem{rem}[thm]{Remark}
\newtheorem{notation}[thm]{Notation}
\newcommand{\N}{\mathcal{N}}
\newcommand{\W}{\mathcal{W}}
\newcommand{\Z}{\mathbb{Z}}
\newcommand{\ov}[1]{\overline{#1}}
\newcommand{\Stab}[2]{\operatorname{Stab}_{#1}(#2)}
\newcommand{\link}{\operatorname{Lk}}
\newcommand{\good}{always restrained}
\title[Short HHG II]{Short hierarchically hyperbolic groups II: quotients and the Hopf property: the broad case of two odd components}
\author{Giorgio Mangioni, Alessandro Sisto}
\date{Source: arXiv:2412.04364v4 (CC BY 4.0)}
\begin{document}
\maketitle
\noindent\emph{Source and licence.} Short hierarchically hyperbolic groups II: quotients and the Hopf property. Version arXiv:2412.04364v4 (CC BY 4.0), distributed under the Creative Commons Attribution 4.0 International licence, \url{http://creativecommons.org/licenses/by/4.0/}, as recorded in the arXiv licence metadata for this exact version and shown on the abstract page \url{https://arxiv.org/abs/2412.04364v4}. Full licence notice: licenses/arxiv-2412.04364-CC-BY-4.0.txt.\par\smallskip

\noindent\textit{Editorial scope.} Counted unit: the article's Proposition with source label \texttt{prop:2\_odd\_broad} (source0.tex lines 2022--2024), the case of two broad odd components in the article's proof that large Artin groups of hyperbolic type are Hopfian, delivered together with the article's complete original proof of it (source0.tex lines 2026--2053, one \texttt{proof} environment of 28 lines, adjacent to the statement, whose case analysis A, B and C closes with the article's completion sentence naming the proposition and the paper's main theorem). No mathematics is written, repaired or re-derived: statement and proof are the article's own text line for line, in the article's own notation (the odd and hanging components, the core, the retraction, the centralisers and dihedrals, the peripheral structure), and no retained line is substituted, re-flowed or omitted. Required context retained uncounted, in source order: the statements of the five results that the counted proposition and its neighbours invoke, namely \texttt{lem:enough} (lines 1563--1566), \texttt{thm:rel\_hyp\_hopf} (lines 1610--1613), \texttt{lem:rel\_hyp\_check} (lines 1661--1664), \texttt{rem:hopfian\_from\_good\_stuff} (lines 1750--1754) and \texttt{thm:XL\_Hopfian} (lines 1850--1853), together with the statements of the two companion cases \texttt{prop:3+\_odd} (lines 1978--1980) and \texttt{prop:2\_odd\_needle} (lines 2009--2011). The proofs of those context results are not delivered. Nothing else of the article is retained: its other sections, its earlier propositions and its figures are omitted. No citation command occurs in any retained line, so the article's bibliography is not delivered and the wrapper carries no bibliography entry, although the article ships one. Every cross-reference inside the retained spans resolves inside the delivered document. Editorial formatting only: an AMS \texttt{amsart} preamble that restates the environments and macros the retained lines use, namely the article's declarations of Theorem, Lemma, Corollary, Proposition, Remark and Notation sharing one counter, the article's macros for the nesting set, the hyperbolic space, the integers, the overline decoration, the stabiliser operator and the link operator, and the article's own macro that expands to the words "always restrained". The article is typeset in its own class with its own fonts and ships style files; no class or style file is shipped with this package and the wrapper is an amsart document. Numbering differs from the article, because the article resets its shared counter in each section and because only a subset of environments is retained: the retained environments print with the article's shared counter in their source order, so the counted proposition prints with the wrapper's own number rather than the article's. No picture environment and no graphical inclusion occurs in any retained line, and the package delivers no image asset. One bundled source file, \texttt{sources/169671/source0.tex}, accompanies the delivered item as the exact source of every retained span. Every retained span is recorded in \texttt{delivery-evidence.json} with method \texttt{exact\_tex}.
\begin{lemma}
\label{lem:enough}
    If $G$ has enough Hopfian quotients then it is Hopfian.
\end{lemma}

\begin{thm}
\label{thm:rel_hyp_hopf}
    Let $G$ be hyperbolic relative to $\Z$-central extensions of hyperbolic groups (including the case that $G$ itself is such an extension). Then $G$ is Hopfian.
\end{thm}

\begin{lemma}
\label{lem:rel_hyp_check}
    If a short HHG $(G, \ov X, \W)$ has discrete product region graph then it is hyperbolic relative to $\{\Stab{G}{v}\}_{v\in V}$, where $V$ is a collection of representatives for some $G$-orbits of vertices with infinite cyclic directions. In particular, if $G$ furthermore has central cyclic directions then it is Hopfian by Theorem~\ref{thm:rel_hyp_hopf}.
\end{lemma}

\begin{rem}\label{rem:hopfian_from_good_stuff}
    Notice that, in the above Lemma, the product region graph of $G/\N$ injects in the graph obtained from $\mathcal{PR}(G)$ after removing the open stars of the vertices corresponding to $v_1,\ldots, v_r$. 
    
    Now, suppose that $G$ has central cyclic directions. In view of the above discussion, if removing all \good{} stabilisers makes $\mathcal{PR}(G)$ discrete, then $G/\N$ is Hopfian by Lemma \ref{lem:rel_hyp_check}. In other words, the quotient $G/\N$ satisfies all requirements of our criterion, Lemma \ref{lem:enough}.
\end{rem}

\begin{thm}\label{thm:XL_Hopfian}
    Let $A_\Gamma$ be an Artin group of large and hyperbolic type, such that every hanging component is either broad or a needle. 
    Then every surjective homomorphism $\phi\colon A_\Gamma\to A_\Gamma$ is an isomorphism.
\end{thm}

\begin{prop}\label{prop:3+_odd}
    Let $A_\Gamma$ be a large Artin group of hyperbolic type, where $\Gamma$ is a connected graph on at least three vertices. Suppose that $\Gamma$ has at least three odd components, and every hanging component is either broad or a needle. Then $A_\Gamma$ is Hopfian.
\end{prop}

\begin{prop}\label{prop:2_odd_needle}
    Let $A_\Gamma$ be a large Artin group of hyperbolic type, where $\Gamma$ is a connected graph on at least three vertices. Suppose that $\Gamma$ has two odd components, one of which is a needle. Then $A_\Gamma$ is Hopfian.
\end{prop}

\begin{prop}\label{prop:2_odd_broad}
    Let $A_\Gamma$ be a large Artin group of hyperbolic type, where $\Gamma$ is a connected graph on at least three vertices. Suppose that $\Gamma$ has two odd components, which are both broad. Then $A_\Gamma$ is Hopfian.
\end{prop}

\begin{proof}
    Let $C=\{a_1,\ldots, a_k\}$ and $C'=\{b_1,\ldots, b_r\}$ be the odd components, let $\phi$ be an epimorphism, and let $g_0\in\ker\phi-\{1\}$. Again, if both $C(a_1)$ and $C(b_1)$ are \good{}, then we can invoke Remark~\ref{rem:hopfian_from_good_stuff} and conclude. So suppose that $C(a_1)$ is not \good{}, so there exists a vertex stabiliser $P$, say with centre generated by $z$, such that $\phi(C(a_1))$ is a virtually $\Z^2$ subgroup of $P$ not contained in any edge group. Arguing as in Proposition~\ref{prop:2_odd_needle}, one gets that $\phi(C)\le P$, and $\phi(A_{a_ib_j})\le P$ for every $i,j$ such that $a_i$ and $b_j$ are joined by an even edge. We now consider three scenarios, A, B, and C, depending on the shape of $\phi(C(b_1))$. 
    \par\medskip
    \textbf{A.}
    Suppose first that $\phi(C(b_1))$ is a $\Z$-central extension of a non-elementary hyperbolic group, so $b_1$ must be sent inside some centre. As standard generators have primitive image in the abelianisation while centres of dihedrals do not, we can assume up to conjugation that $\phi(b_1)=b_1$ or $\phi(b_1)=a_1$. In the former case, $\phi$ preserves the normal closure of $C'$, so it induces a self-map $\ov\phi$ of the quotient $A_C$, obtained by retracting onto $C$. Then one can run the proof of Proposition~\ref{prop:3+_odd} verbatim, to get that $A_\Gamma$ has “enough” \good{} vertex stabilisers. 
\par\medskip

    In the latter case, $\phi^2(C(b_1))\le \phi(C(a_1))\le P$. It now matters what $P$ is, taking into account that it cannot be an odd dihedral as $\phi(C(a_1))$ must map to $\Z^2$ in the abelianisation. If $P$ is a conjugate of $C(b_1)$ then we can argue as above, with $\phi^2$ replacing $\phi$. If instead $P$ is conjugated to $Q\in\{C(a_1),A_{a_ib_j}\}$ by some $g\in G$, up to composing $\phi$ with the conjugation by $g^{-1}$ we can assume that $\phi(Q)\le Q$. Then we have that:
    \begin{itemize}
        \item $\phi^2(C)\le \phi(Q)\le Q$;
        \item $\phi^2(C')\le Q$, because $\phi^2(C(b_1))\le \phi(C(a_1))$ is contained in a $\Z^2$ subgroup which is not an edge group. 
    \end{itemize}
    Hence $\phi^2(A_\Gamma)\le Q$, violating surjectivity.

\par\medskip
    \textbf{B.}
    Suppose now that $\phi(C(b_1))$ is virtually $\Z^2$, but not contained in an edge group. Since $\phi(C(b_1))$ contains $\phi(z_{a_1b_1})$, which belongs to the non-edge $\Z^2$ subgroup $\phi(C(a_1))$ of $P$, we must have that $\phi(C(b_1))\le P$. Then again, using that $\phi(C(b_1))$ is not contained in an edge group, we get that $\phi(C')\le P$, which combined with $\phi(C)\le P$ violates surjectivity.
    \par\medskip
    \textbf{C.}
    Suppose finally that $\phi(C(b_1))$ is a finite-index subgroup of an edge group. As above, since $\phi(C(b_1))$ contains $\phi(z_{a_1b_1})$ we must have that $\phi(C(b_1))\le P$. We now consider the possible conjugacy types of $P$.
    \begin{itemize}
        \item Suppose first that $P$ is conjugate to $C(a_1)$, and up to composing $\phi$ with an inner automorphism we can indeed assume that $P=C(a_1)$. Then $\phi^2(C(b))\le \phi(C(a))\le P$ is a non-edge, virtually $\Z^2$ subgroup, and again this implies that $\phi(C')\le P$, contradicting surjectivity.
        \item Suppose now that $P$ is conjugate to some dihedral, which must be of the form $A_{a_ib_j}$ (again because any other dihedral has cyclic image in $A_\Gamma^{ab}$). Pick $b_2\in C'$ which is connected to $b_1$ by an odd edge. Then $\phi(z_{b_1b_2})\in P$, and as parabolics are root closed we must have that $\phi(b_1b_2)\in P$. Hence $\phi(C(b_2))$, which is obtained by conjugating $\phi(C(b_1))$ by a power of $\phi(b_1b_2)$, is also nested in the subgroup $P$. Proceeding this way, we eventually get that $\phi(C')\le P$, and once more $\phi$ could not be surjective. 
        \item The only case left is when $P$ is conjugate to $C(b_1)$, and again we can indeed assume that $P=C(b_1)$ up to composing $\phi$ with a conjugation. Say $C(b_1)=\Stab{A_\Gamma}{v}$ for some $v\in \ov X^{(0)}$. As $\phi(C(b_1))$ lies in some edge group, there must be some $w\in\link_{\ov X}(v)$ such that, if we set $Q=\Stab{A_\Gamma}{w}$, then $\phi(C(b_1))\le C(b_1)\cap Q$. Since $\ov X$ is bipartite, $Q$ must be conjugate to a dihedral, so the same trick as above shows that $\phi(C')\le Q$. If we show that $\phi(Q)\le C(b_1)$ then $\phi^2(C')\le C(b_1)$, and $\phi^2(C)\le \phi(C(b_1))\le C(b_1)$. This would then again contradict surjectivity. 
        \\
        To prove that $\phi(Q)\le C(b_1)$, we first notice that the image of $Q$ in the abelianisation must have rank 2, or it could not contain $\phi(C(b_1))$; hence there exists $i,j$ such that $Q$ is conjugate to $A_{a_ib_j}$. In turn, since $\phi(A_{a_ib_j})\le C(b_1)$, there must be some $P'$, which is a conjugate of $C(b_1)$, such that $\phi(Q)\le P'$. But $\phi(Q)$ contains $\phi^2(C(b_1))$, which is virtually $\Z^2$ and lies inside $\phi(C(b_1))\le C(b_1)$. As $\ov X$ is bipartite, any two different conjugates of $C(b_1)$ intersect along a virtually cyclic subgroup, so we must have that $P'=C(b_1)$, as required.
    \end{itemize}
This concludes the proof of Proposition~\ref{prop:2_odd_broad}, and in turn of Theorem~\ref{thm:XL_Hopfian}. \end{proof}

\end{document}
